\documentclass[11pt]{article}
\usepackage[margin=1in]{geometry}
\usepackage{amsmath,amssymb,mathtools}
\usepackage{bm}
\usepackage{booktabs}

\title{Detailed Writeup of VI and MF-VI in ROMTools\\
\large Focus on Newton Updates in Natural Coordinates}
\author{}
\date{}

\begin{document}
\maketitle

\section{Scope and Notation}
This note describes the two inverse workflows implemented in this module:
\begin{itemize}
  \item \texttt{run\_vi}: single-fidelity stochastic variational inference.
  \item \texttt{run\_mf\_vi}: multi-fidelity variational inference with MFMC-style control variates.
\end{itemize}
The emphasis is on the Newton update with \texttt{newton\_metric = "natural"}.

Let:
\begin{itemize}
  \item \(y \in \mathbb{R}^{n_q}\): observation vector.
  \item \(f(\theta)\in\mathbb{R}^{n_q}\): model QoI at parameter \(\theta\).
  \item \(e(\theta)=y-f(\theta)\): model error.
  \item \(\Sigma_y\): observation covariance.
  \item \(\lambda_{\mathrm{cov}}\): covariance regularization.
  \item \(P = (\Sigma_y + \lambda_{\mathrm{cov}} I)^{-1}\) (pseudo-inverse in implementation).
\end{itemize}

For each sample \(\theta^{(m)}\), the implemented log-likelihood is
\begin{equation}
\log p(y\mid \theta^{(m)}) = -\frac12\, e(\theta^{(m)})^\top P\, e(\theta^{(m)}).
\end{equation}

\section{Background: Variational Inference and REINFORCE}
\subsection{Why Variational Inference}
Given prior \(p(\theta)\) and likelihood \(p(y\mid\theta)\), Bayes' rule gives
\begin{equation}
p(\theta\mid y)=\frac{p(y\mid\theta)p(\theta)}{p(y)}.
\end{equation}
For nonlinear QoI models, this posterior is typically intractable to normalize or sample repeatedly.
Variational inference replaces posterior sampling with optimization:
choose a tractable family \(q(\theta;\lambda)\), then optimize \(\lambda\) so
\(q\) is close to \(p(\theta\mid y)\).

\subsection{ELBO Viewpoint}
Using \(\mathrm{KL}(q\|p)\ge 0\), maximizing ELBO is equivalent to minimizing a KL divergence:
\begin{align}
\log p(y)
&= \mathcal{L}(\lambda) + \mathrm{KL}\!\left(q(\theta;\lambda)\,\|\,p(\theta\mid y)\right),\\
\mathcal{L}(\lambda)
&= \mathbb{E}_{q}\!\left[\log p(y,\theta)-\log q(\theta;\lambda)\right].
\end{align}
In this implementation, the prior term is not modeled explicitly in the objective; the optimized quantity is
\(s\,\mathbb{E}_{q}[\log p(y\mid\theta)] + s\,\mathcal{H}[q]\), where
\(\mathcal{H}[q]=-\mathbb{E}_q[\log q]\).

\subsection{Why Score-Function (REINFORCE) Gradients}
For an objective \(\mathbb{E}_{q_\lambda}[f(\theta)]\), the score-function identity gives
\begin{equation}
\nabla_\lambda \mathbb{E}_{q_\lambda}[f(\theta)]
= \mathbb{E}_{q_\lambda}\!\left[f(\theta)\,\nabla_\lambda \log q_\lambda(\theta)\right].
\end{equation}
This is the REINFORCE estimator. It is useful here because:
\begin{itemize}
  \item The forward model is treated as a black-box QoI map.
  \item Gradients of \(f(\theta)\) w.r.t. \(\theta\) are not required.
  \item The same machinery applies to single- and multi-fidelity objectives.
\end{itemize}

\subsection{Variance Reduction in REINFORCE}
The implementation uses standard control variates:
\begin{equation}
\mathbb{E}\!\left[(f-b)\,\nabla_\lambda \log q\right]
=
\mathbb{E}\!\left[f\,\nabla_\lambda \log q\right],
\end{equation}
because \(\mathbb{E}_q[\nabla_\lambda \log q]=0\).
This reduces variance without bias.
Available choices are \texttt{none}, \texttt{loo}, and \texttt{optimal}; MF-VI adds MFMC control variates across fidelity levels.

\subsection{Second-Order Extension}
Newton updates require curvature information in parameter space.
This code builds a score-function Hessian estimator using \(\nabla \log q\) and
\(\nabla^2 \log q\), then stabilizes the solve with Hessian projection/regularization.
Natural-coordinate Newton further rescales this system by Fisher-inspired factors.

\section{Variational Family and Parameterization}
The variational parameters are \((\mu,\ell)\), where \(\ell=\log \sigma\), \(\sigma=\exp(\ell)\), and \(\sigma_i\in[\sigma_{\min},\sigma_{\max}]\) by clipping.

\subsection{Diagonal Gaussian}
\begin{equation}
q(\theta;\mu,\ell)=\mathcal{N}\!\left(\mu,\operatorname{diag}(\sigma^2)\right),\qquad
\theta=\mu+\sigma\odot z,\; z\sim\mathcal{N}(0,I).
\end{equation}

\subsection{Multivariate Gaussian with Fixed Correlation}
The implementation supports a fixed correlation Cholesky factor \(L\), with
correlation matrix \(R=LL^\top\):
\begin{equation}
\theta=\mu + D_\sigma L z,\qquad z\sim\mathcal{N}(0,I),\qquad
D_\sigma=\operatorname{diag}(\sigma).
\end{equation}
Only \((\mu,\ell)\) are optimized; \(L\) is treated as fixed during optimization.

\section{ELBO Used in Both Workflows}
With scaling factor \(s=\texttt{elbo\_scaling\_factor}\), the optimized objective is
\begin{equation}
\mathcal{J}(\mu,\ell) = s\,\mathbb{E}_{q}[\log p(y\mid\theta)] + s\,\mathcal{H}[q].
\end{equation}
Entropy terms:
\begin{align}
\mathcal{H}_{\text{diag}}[q] &= \sum_i \ell_i + \frac{d}{2}\left(1+\log(2\pi)\right),\\
\mathcal{H}_{\text{mv}}[q] &= \sum_i \ell_i + \log\det L + \frac{d}{2}\left(1+\log(2\pi)\right).
\end{align}
The implementation multiplies entropy by the same \(s\).

\section{Single-Fidelity VI: Gradient and Hessian Estimators}
\subsection{Score Functions}
Define normalized coordinates \(u=(\theta-\mu)\oslash \sigma\).

\paragraph{Diagonal family}
\begin{align}
\nabla_\mu \log q &= (\theta-\mu)\oslash \sigma^2,\\
\nabla_\ell \log q &= u\odot u - \mathbf{1}.
\end{align}

\paragraph{Fixed-correlation multivariate family}
Let \(C=R^{-1}\), \(v=C\,u\). Then:
\begin{align}
\nabla_\mu \log q &= v \oslash \sigma,\\
\nabla_\ell \log q &= u\odot v - \mathbf{1}.
\end{align}

\subsection{Baseline Handling}
Three baseline choices are implemented for variance reduction:
\begin{itemize}
  \item \texttt{none}: no baseline.
  \item \texttt{loo}: leave-one-out scalar baseline per sample.
  \item \texttt{optimal}: componentwise optimal baseline \(b=\frac{\mathbb{E}[w s^2]}{\mathbb{E}[s^2]}\) for each score component.
\end{itemize}
Here \(w = s\log p(y\mid\theta)\).

\subsection{Gradient Estimator}
For each block (\(\mu\), \(\ell\)):
\begin{equation}
g = \frac{1}{N}\sum_{m=1}^{N} (w_m-b_m)\,s_m.
\end{equation}
For the \(\ell\)-block, entropy contributes \(+s\mathbf{1}\), yielding
\begin{equation}
g_\ell \leftarrow g_\ell + s\mathbf{1}.
\end{equation}

\subsection{Hessian Estimator (Score-Function Form)}
The implementation builds the block Hessian via
\begin{equation}
H \approx \frac1N\sum_{m=1}^{N}
(w_m-b_m)\left(s_m s_m^\top + \nabla^2 \log q_m\right),
\end{equation}
assembled into \((\mu,\ell)\) blocks:
\[
H=
\begin{bmatrix}
H_{\mu\mu} & H_{\mu\ell}\\
H_{\ell\mu} & H_{\ell\ell}
\end{bmatrix}.
\]

\paragraph{Diagonal-family second derivatives}
\begin{align}
\frac{\partial^2 \log q}{\partial \mu_i\partial\mu_k}
&= -\delta_{ik}\sigma_i^{-2},\\
\frac{\partial^2 \log q}{\partial \ell_i\partial\ell_k}
&= -2\delta_{ik}u_i^2,\\
\frac{\partial^2 \log q}{\partial \mu_i\partial\ell_k}
&= -2\delta_{ik}\frac{\partial \log q}{\partial\mu_i}.
\end{align}

\paragraph{Fixed-correlation second derivatives}
With \(C=R^{-1}\), \(u=(\theta-\mu)\oslash \sigma\), \(v=Cu\):
\begin{align}
\frac{\partial^2 \log q}{\partial \mu_i\partial\mu_k}
&= -\frac{C_{ik}}{\sigma_i\sigma_k},\\
\frac{\partial^2 \log q}{\partial \ell_i\partial\ell_k}
&= -u_i u_k C_{ik} - \delta_{ik}(u_i v_i),\\
\frac{\partial^2 \log q}{\partial \mu_i\partial\ell_k}
&= -\frac{u_k C_{ik}}{\sigma_i} - \delta_{ik}\frac{v_i}{\sigma_i}.
\end{align}
These are exactly the terms used in the new multivariate-Newton implementation.

\section{Newton Step and Natural Coordinates}
\subsection{Standard Newton Path}
The optimization vector is
\begin{equation}
\vartheta=\begin{bmatrix}\mu\\\ell\end{bmatrix},\qquad
g=\begin{bmatrix}g_\mu\\g_\ell\end{bmatrix},\qquad
H=\nabla_\vartheta^2 \mathcal{J}.
\end{equation}
The solver computes \(\Delta\) from \(H\Delta=g\), then tests \(\vartheta+\alpha\Delta\).

\subsection{Hessian Projection/Regularization}
Before solving, \(H\) is projected to a numerically safe SPD operator:
\begin{itemize}
  \item diagonal case: \(H_{ii}\leftarrow \max(|H_{ii}|,\lambda_{\mathrm{newton}})\),
  \item dense case: symmetrize, eigendecompose, clamp \(|\lambda_j|\ge \lambda_{\mathrm{newton}}\), reconstruct.
\end{itemize}
Then solve the linear system.

\subsection{Natural-Coordinate Newton (\texttt{newton\_metric = "natural"})}
The code uses scale
\begin{equation}
S=\operatorname{diag}\!\big(\sigma_1,\dots,\sigma_d,\sqrt{1/2},\dots,\sqrt{1/2}\big),
\end{equation}
which is \(F^{-1/2}\)-consistent for diagonal Gaussian Fisher blocks in \((\mu,\ell)\).

Newton is performed in scaled coordinates:
\begin{align}
\tilde g &= S g,\\
\tilde H &= S H S,\\
\tilde H\,\tilde\Delta &= \tilde g,\\
\Delta &= S\tilde\Delta.
\end{align}
This is an approximate natural-Newton preconditioning strategy (without full connection terms).

\section{Step Proposal and Stochastic Line Search}
\subsection{Candidate Update}
Given direction \((d_\mu,d_\ell)\) and step size \(\alpha\):
\begin{align}
\mu^+ &= \mu + \alpha d_\mu,\\
\ell^+ &= \ell + \operatorname{clip}\!\left(\alpha\,\eta_\ell\, d_\ell,\,-\Delta_{\ell,\max},\Delta_{\ell,\max}\right),
\end{align}
with \(\eta_\ell=\texttt{log\_std\_learning\_rate\_factor}\) in \texttt{run\_vi}, and \(\eta_\ell=1\) in \texttt{run\_mf\_vi}.

Predicted slope used by Armijo logic:
\begin{equation}
\hat s = g_\mu^\top d_\mu + \eta_\ell\, g_\ell^\top d_\ell
\quad(\eta_\ell=1\text{ in MF-VI}).
\end{equation}

\subsection{Acceptance Rules}
\paragraph{Legacy}
\begin{equation}
\mathcal{J}^+ \ge \mathcal{J} - (\rho-1)|\mathcal{J}|,
\end{equation}
where \(\rho=\texttt{relaxation\_parameter}\).

\paragraph{Stochastic nonmonotone Armijo}
Let \(\mathcal{J}_{\mathrm{ref}}=\max\) over a recent accepted window.
The target is
\begin{equation}
\mathcal{J}_{\mathrm{target}}
= \mathcal{J}_{\mathrm{ref}}
+ c\,\alpha\,\max(\hat s,0)
- \kappa\,\Delta_{\mathrm{SE}},
\end{equation}
with \(c=\texttt{line\_search\_armijo\_coefficient}\),
\(\kappa=\texttt{line\_search\_uncertainty\_sigma}\).
Accept if \(\mathcal{J}^+\ge \mathcal{J}_{\mathrm{target}}\).

Step-size adaptation:
\begin{itemize}
  \item accept: \(\alpha\leftarrow \min(\gamma_{\uparrow}\alpha,\alpha_{\max})\),
  \item reject: \(\alpha\leftarrow \alpha/\gamma_{\downarrow}\).
\end{itemize}
Sample size can grow after repeated failures via \texttt{line\_search\_sample\_growth\_factor}.

\section{MF-VI: Multi-Fidelity Construction}
\subsection{Three Sample Sets per Iteration}
At current variational parameters:
\begin{itemize}
  \item FOM set (\(N_H\)): high-fidelity samples.
  \item ROM-base set (\(N_H\)): low-fidelity samples paired in count with FOM.
  \item ROM-extra set (\(N_L\)): additional low-fidelity samples.
\end{itemize}
ROM-base can be \texttt{coupled} (reuse FOM points) or \texttt{separate} (independent draw).

\subsection{ROM Rebuild Logic}
The ROM is rebuilt from recent FOM training data when ROM mean-QoI relative error exceeds \texttt{rom\_tolerance}.
This couples ROM management with the VI iteration.

\subsection{MFMC Estimator for Vector/Tensor Quantities}
For any per-sample term \(T\) (vector or matrix valued), with high term \(T_H\), low-base \(T_{L,b}\), low-extra \(T_{L,e}\), define
\begin{equation}
\widehat{\mathbb{E}}[T]
= \overline{T_H}
 + \alpha_T\left(\overline{T_{L,\mathrm{full}}}-\overline{T_{L,b}}\right),
\qquad
T_{L,\mathrm{full}}=[T_{L,b};T_{L,e}].
\end{equation}
Componentwise control variate coefficient:
\begin{equation}
\alpha_T = \frac{\operatorname{Cov}(T_H,T_{L,b})}{\operatorname{Var}(T_{L,b})},
\end{equation}
with safe fallback when \(\operatorname{Var}(T_{L,b})\) is too small.
If control variates are disabled, \(\alpha_T=1\). If \(N_L=0\), estimator reduces to \(\overline{T_H}\).

\subsection{MF-VI Gradient/Hessian}
The above MFMC estimator is applied to:
\begin{itemize}
  \item gradient mean/log-std terms,
  \item Hessian \(\mu\mu\), \(\ell\ell\), and \(\mu\ell\) block terms.
\end{itemize}
Entropy contribution \(+s\mathbf{1}\) is still added to \(g_\ell\).

\subsection{MF-VI ELBO and Uncertainty Term}
When \(N_L>0\), mean log-likelihood estimate is
\begin{equation}
\widehat{\mathbb{E}}[\log p]
= \overline{\log p_H}
 + \overline{\log p_{L,\mathrm{full}}}
 - \overline{\log p_{L,b}}.
\end{equation}
ELBO estimate is \(s\,\widehat{\mathbb{E}}[\log p] + s\,\mathcal{H}[q]\).

For stochastic line-search uncertainty in MF-VI, standard errors are combined as
\begin{equation}
\mathrm{SE}_{\mathrm{mfmc}}
= \sqrt{\mathrm{SE}_H^2 + \mathrm{SE}_{L,\mathrm{full}}^2 + \mathrm{SE}_{L,b}^2},
\end{equation}
and then scaled by \(|s|\).

\section{Natural-Coordinate Newton in MF-VI}
The Newton solver and natural-coordinate scaling are exactly the same as in VI:
\begin{equation}
S=\operatorname{diag}(\sigma,\sqrt{1/2}\mathbf{1}),\quad
\tilde g= Sg,\quad \tilde H=SHS,\quad \Delta=S\tilde H^{-1}\tilde g,
\end{equation}
but now \(g,H\) are MFMC estimates built from FOM/ROM sample sets.

\section{Implementation-Oriented Summary}
\begin{itemize}
  \item Both VI and MF-VI optimize \((\mu,\ell)\) with score-function gradients and Hessians.
  \item Both now support Newton with diagonal and fixed-correlation multivariate families.
  \item Natural-coordinate Newton is implemented as Fisher-inspired scaling of gradient/Hessian before solve.
  \item Stochastic nonmonotone Armijo line search stabilizes noisy objective evaluations.
  \item MF-VI reduces estimator variance/cost via hierarchical low-fidelity corrections and online ROM management.
\end{itemize}

\section{Configuration Snippet}
\begin{verbatim}
optimizer_method = "newton"
optimizer_config = VINewtonOptimizerConfig(
    newton_metric="natural",
    newton_regularization=1e-1,
)
line_search_method = "stochastic_nonmonotone"
\end{verbatim}
For MF-VI, use \texttt{run\_mf\_vi(...)} with the same optimizer config plus ROM options.

\end{document}
